\documentclass[12pt, a4paper]{article}
% --- Packages -----------------------------------------------------------------
\usepackage[T1]{fontenc}
\usepackage[utf8]{inputenc}
\usepackage[french, main=english]{babel}
\usepackage[autostyle=true]{csquotes}
% Tipografía: Tipografíca de Visual Complex Analysis (más o menos)
\usepackage{amsmath, amsthm, amssymb}
\usepackage{mathpazo}
\usepackage{eulerpx}
% Tipografía estándar: \usepackage{amsmath, amsthm, amssymb, amsfonts}
\usepackage[hidelinks]{hyperref}

% --- Document Information -----------------------------------------------------
\title{On the Iteration of Rational Substitutions and Poincaré Functions}
\author{ Original Note by S. Lattès (1918) \thanks{ Originally published in
    French as \emph{``Sur l'itération des substitutions rationnelles et les
    fonctions de Poincaré''}, \emph{Comptes Rendus de l'Académie des Sciences},
    Paris, Volume 166, Pages 26--28 (January 7, 1918). } }
\date{January 7, 1918}

% --- Document -----------------------------------------------------------------
\begin{document}

\subsection*{Mathematical Analysis. --- On the Iteration of Rational
Substitutions and Poincaré Functions. Note by M. S. Lattés.}

Let $z_1 = R(z)$ be a rational map with distinct fixed points; let $\alpha$ be a
fixed point with multiplier S, where $|S| > 1$. There is always at least one
such point\footnote{\foreignlanguage{french}{Cf. Fatou, \emph{Sur les
substitutions rationnelles} (Comptes rendus, t. 165, 1917, p. 993).}}. There
then exists a function $\theta(u)$ which is meromorphic [or entire if $R(z)$ is
a polynomial] satisfying, for all $u$, the functional equation:
\begin{equation}
    \label{eq:1}
    \theta(Su) = R[\theta(u)]
\end{equation}
and taking the value $\alpha$ at $u = 0$. The existence of $\theta(u)$, which we
shall call the \emph{Poincaré function} relative to the invariant point
$\alpha$, was established by
Poincaré\footnote{\foreignlanguage{french}{Poincaré, \emph{Sur une classe
nouvelle de transcendantes uniformes} (Journal de Mathématiques pures et
appliquées, 1890).}}.

\vspace{1em}
{
\small
Poincaré, however, imposes on the desired function $\theta(u)$ the condition
$\theta'(0) \neq 0$. If $\theta'(0)$ is zero and the first non-zero derivative
for $u = 0$ is of order $p$, one still finds a meromorphic function $\theta(u)$
satisfying equation (\ref{eq:1}), but the number $S$ appearing in this equation---which
can still be called the multiplier---is related to $R'(\alpha)$ by the relation
$S^p = R'(\alpha)$; the function $\theta(u)$ is then a meromorphic or entire
function of $u^p$. For example, if the given substitution is $z_1 = 2z^2 - 1$,
we can set $z = \cos u$, yielding $z_1 = \cos 2u$; here $\theta(u)$ is an entire
function of $u^2$ [since $\cos u = 1 - \frac{u^2}{2} + \dots$] and we have
$R'(\alpha) = 4 = S^2$ [where $\alpha = 1$, $R'(1) = 4$, and the multiplier is
$S = 2$].
}
\vspace{1em}

The Poincaré function $\theta(u)$ is nothing other than the inverse function of
the Schröder function $\varphi(z)$, whose existence in the neighborhood of the
point $\alpha$ was established by Mr. K\oe~nigs and which satisfies the
equation:
\[
    \varphi[R(z)] = S\varphi(z).
\]
But, while the Poincaré function, when analytically continued from $u = 0$,
gives rise to a single-valued (uniform) meromorphic or entire transcendental
function, its inverse, the Schröder function, continued from $z = \alpha$, is by
that very fact a multivalued (multiform) analytic function. Hence, there is an
advantage to substituting the Poincaré function for the Schröder function if one
wishes to study iteration starting from an arbitrary initial point $z$ in the
plane.

\newpage

For each value of $z$, the relation $z = \theta(u)$ yields at least one value of
$u$, except for at most two exceptional values of $z$, according to Picard's
Theorem. We then obtain, in parametric form, the values for the point $z$ and
its successive iterates [consequents] $z_n$:
\begin{equation*}
    z   = \theta(u), \quad z_1 = \theta(Su), \quad z_2 = \theta(S^2 u), \quad \ldots, \quad z_n = \theta(S^n u).
\end{equation*}
These same formulas, for negative integer $n$, provide a sequence of successive
antecedents [preimages] $z_{-n}$ of the point $z$ which, as $n \to \infty$, have
a limit equal to $\theta(0)$, which is $\alpha$. Thus, \emph{one can choose, from
among the possible successive preimages of any given point, a sequence of
successive preimages $z_{-n}$ tending, as $n \to \infty$, toward any of the
invariant points $\alpha$ whose multiplier has a modulus greater than $1$}.

The possible exceptional points $\alpha, \beta$ of the function $\theta(u)$ are
either invariant points of the substitution\footnote{For example, for the
substitution $z_1 = \frac{2z}{1-z^2}$, the Poincaré function is $z = \tan u$,
with multiplier $2$; this function admits the exceptional values $i$ and $-i$,
which are invariant points of the substitution.} or two points forming a
periodic cycle of order 2.

\emph{The problems relating to the iteration of a given substitution are thus
transformed into problems concerning the growth of the function $\theta(u)$}. To
obtain categories of rational substitutions whose iteration is easy to study, it
is convenient to choose substitutions that admit a Poincaré function belonging
to the most common meromorphic (or entire) functions. Thus, if we assume $S$ to
be an integer, and $\theta(u)$ to be one of the functions $\tan u$, $\cos u$,
$\wp(u)$ [the Weierstrass elliptic function], or a homographic transformation
[Möbius transformation] of these functions, we obtain rational substitutions for
which the method of parametric iteration allows us to completely resolve the
fundamental problem of iteration. This problem, in its general form, can be
stated as follows:

\vspace{1em}
\emph{Determine the limit set $E'$ (derived set) of the set $E$ of iterates
$z_n$ of an arbitrarily given point $z$.}
\vspace{1em}

The set $E'$ contains the iterates of its various points; on the other hand, we
know that $E'$, being a closed set, is the sum of a perfect set $E_1'$ and a
countable set $E_2'$. The set $E_1'$ also contains the iterates of its various
points, while for points in $E_2'$, their iterates may belong to $E_1'$. In the
various examples listed above, we are able to completely determine $E_1'$ and
$E_2'$, and determine what conditions $z$ must satisfy for either of the sets
$E_1'$ or $E_2'$ to be empty, or for a given point in the plane to belong to
$E'$. These various conditions depend on the arithmetic properties of the number
$z$.

\vspace{1em}
{
\small
Consider, for example, the substitution:
\begin{equation}
    \label{eq:2}
    z_1 = \frac{(z^2+1)^2}{4z(z^2-1)}
\end{equation}
which is obtained by setting:
\[
    z = \wp(u), \quad z_1 = \wp(2u) \quad \text{with} \quad g_2 = 4, \ g_3 = 0.
\]
If $2\omega$ denotes the real period, the other period is $2i\omega$. If we set:
\[
    z = \wp(u) = \wp(2\omega v + 2i\omega w),
\]
where $v$ and $w$ are real, then to determine $E'$, we must know the
representations of $v$ and $w$ in the binary numeral system (base 2). The set
$E_1'$ is formed by drawing the orthogonal grid of curves $v = \text{constant}$
and $w = \text{constant}$ (Cartesian ovals) and excluding from the plane the
interior points of a countable infinity of quadrilaterals---contiguous or
not---bounded by the curves of this grid. The common boundary between two
excluded contiguous quadrilaterals must generally also be excluded. It is for
the determination of the arcs of the ovals that bound the excluded domains that
we must know the binary representations of $v$ and $w$. We can thus completely
solve the various iteration problems relating to (\ref{eq:2}).
}
\vspace{1em}

The method of parametric iteration, using generalized Poincaré functions, can be
applied in very broad cases to rational substitutions of two variables.

\end{document}